% Clean, formal report-style document class and packages
\documentclass[11pt, twoside, a4paper]{article}
\usepackage[utf8]{inputenc}
\usepackage[a4paper, margin=2.5cm]{geometry}

% Essential packages for a formal report
\usepackage{mathtools, amsmath, amssymb, derivative}    % Math environments
\usepackage{float, graphicx, booktabs, subcaption, array} % Figures and tables
\usepackage{listings}                                    % Source code
\usepackage{xcolor}                                     % Colors
\usepackage{titlesec}                                   % Section formatting
\usepackage{pdfpages}                                   % PDF inclusion
\usepackage[style=ieee]{biblatex}                        % Bibliography
\usepackage[hidelinks]{hyperref}                        % Links
\usepackage{siunitx}                                    % SI units
\usepackage[version=4]{mhchem}                          % Chemical formulae
\usepackage{enumitem}                                   % List environments
\usepackage{menukeys}                                   % Menu formatting
\usepackage{tcolorbox}                                  % Boxed environments
% \usepackage{caption}
% \captionsetup[table]{position=top}

% Bibliography resources
\addbibresource{references.bib}
\addbibresource{sample.bib}

% Color definitions for code listings
\definecolor{codegreen}{rgb}{0,0.6,0}
\definecolor{codegray}{rgb}{0.5,0.5,0.5}
\definecolor{codepurple}{rgb}{0.58,0,0.82}
\definecolor{backcolour}{rgb}{1,1,1}

% Code listing style
\lstdefinestyle{mystyle}{
    backgroundcolor=\color{backcolour},   
    commentstyle=\color{codegreen},
    keywordstyle=\color{magenta},
    numberstyle=\tiny\color{codegray},
    stringstyle=\color{codepurple},
    basicstyle=\ttfamily\footnotesize,
    breakatwhitespace=false,         
    breaklines=true,                 
    captionpos=b,                    
    keepspaces=true,                 
    numbers=left,                    
    numbersep=5pt,                  
    showspaces=false,                
    showstringspaces=false,
    showtabs=false,                  
    tabsize=2,
    frame=topbottom
}
\lstset{style=mystyle}

% Custom box for outputs
\newtcolorbox{outputbox}[1][]{
    colback=white,
    title=Output,
    #1
}

% Math operators
\DeclareMathOperator{\proj}{proj}

% Custom column types
\newcommand*{\carry}[1][1]{\overset{#1}}
\newcolumntype{B}[1]{r*{#1}{@{\,}r}}

% Formal report formatting
\setlength{\parindent}{0pt}                             % No paragraph indentation
\setlength{\parskip}{0.8em}                             % Space between paragraphs

% Section spacing for formal reports
\titlespacing*{\section}{0pt}{2.5ex plus 1ex minus 0.2ex}{1.3ex plus 0.2ex}
\titlespacing*{\subsection}{0pt}{2.25ex plus 1ex minus 0.2ex}{1ex plus 0.2ex}
\titlespacing*{\subsubsection}{0pt}{1.5ex plus 1ex minus 0.2ex}{0.75ex plus 0.1ex}

% List formatting for formal reports
\setlist[itemize]{leftmargin=*, itemsep=0.5em, topsep=0.5em}
\setlist[enumerate]{leftmargin=*, itemsep=0.5em, topsep=0.5em}
\setlist[itemize]{label=\textbullet}                    % Use bullets for lists
\setlist[itemize,2]{label=\textendash}                  % En-dash for second level

% Hide underfull warnings
\raggedbottom

% Custom theorem environments
\newtheorem{definition}{Definition}

% Custom commands for document information
\newcommand{\TitleName}{Business Plan and Engineering Design Progress Report}
\newcommand{\unitTime}{Semester 1, 2026}
\newcommand{\unitName}{EGH419: Advanced Design and Entrepreneurship}
\newcommand{\documentAuthors}{Xanthea Vazey: n11525762}

\begin{document}
% -------------------------------------------------
%  Title Page
% -------------------------------------------------
\begin{titlepage}
    \thispagestyle{empty}
    \vspace*{0.2\textheight}
    
    \noindent
    \begin{minipage}[t]{0.03\textwidth}
        \rule[-25ex]{1pt}{140pt}   % longer vertical line
    \end{minipage}
    \hspace{1em}
    \begin{minipage}[t]{0.9\textwidth}
        \raggedright
        \normalsize{\unitTime} \\[2ex]
        \Huge\textbf{\TitleName} \\[2ex]
        \normalsize{\unitName}
    \end{minipage}
    
    \vfill
    \begin{flushleft}
        \documentAuthors
    \end{flushleft}
\end{titlepage}

% Page numbering setup
\pagenumbering{roman}
\newpage
\pagenumbering{arabic}

% -------------------------------------------------
%  Main Content
% -------------------------------------------------


\section{Product Name}

The proposed product developed by Group 40 is titled \textbf{PowerStake}. PowerStake is a smart power management system designed to help households reduce electricity costs through appliance-level monitoring, automated outlet control, predictive insights, and a user-friendly mobile dashboard. The concept combines affordable hardware with intelligent software to reduce standby energy waste, improve convenience, and identify potential appliance faults before failure occurs.

The Group 40 project team is structured using functional roles to allocate responsibilities across business, technical, and design tasks. Team members and primary subsystem responsibilities are listed below:

\begin{itemize}
    \item \textbf{Handler - Erin Keur:} Project coordination and power usage reduction strategy.
    
    \item \textbf{Hustler - Ethan Bailey:} Commercial strategy, customer value proposition, and mobile dashboard development.
    
    \item \textbf{Hacker - Andrew Pollard:} Hardware development including power measurement and supply systems.
    
    \item \textbf{Hipster - Xanthea Vazey:} User experience, industrial design, and automatic power-off scheduling functions.
    
    \item \textbf{Helper - Jeremy Khin Zaw:} Cross-functional support across all project areas, with primary responsibility for fault detection systems.
\end{itemize}

This team structure enables parallel development of the core PowerStake subsystems while maintaining shared ownership of prototype integration and validation activities.

% \subsection{Product Overview}

% PowerStake is a smart household energy management system designed to reduce unnecessary electricity costs through appliance-level monitoring, automated control, and user-friendly 
% recommendations. The primary problem addressed is rising household electricity costs caused by standby power consumption, inefficient appliance usage, and a lack of actionable 
% energy insights. Many users understand that electricity is expensive but do not have the time or motivation to manually switch devices off each day  \cite{smartplugmarket}. The intended end-user is a 
% cost-conscious household occupant or individual seeking lower bills through simple automation rather than technical energy management. 

\section{Target TRL}

The current readiness level of the Automatic Power-off Scheduling subsystem is estimated at TRL 2. At present, the concept has been defined and supported through 
preliminary research into existing technologies such as smart plugs, programmable timers, relay modules, and mobile-controlled outlet systems. However, no integrated 
prototype has yet been constructed or experimentally validated by the project team.

Given the remaining project schedule, the realistic target for this assessment is TRL 3 progressing toward early TRL 4. This would involve developing and 
demonstrating a proof-of-concept prototype using commercially available microcontrollers (e.g. ESP32), relay modules, and basic power sensing inputs. Initial 
validation may be performed on low-voltage test circuits or simulated appliance loads rather than direct 240 V mains systems, reducing safety risk and development 
complexity.

A successful prototype outcome would demonstrate the core subsystem functions of detecting inactive loads, applying simple shutoff logic, manually restoring power,
 and communicating with a mobile dashboard interface. Supporting work may also include enclosure concepts developed through CAD modelling and early physical packaging 
 studies for future iterations.

Achieving a full TRL 4 mains-ready prototype within the current timeframe is considered ambitious due to hardware procurement, electrical safety constraints, and 
software integration requirements. A more practical near-term pathway would be completion of a validated bench-scale prototype during this unit, followed by future 
progression to TRL 4-5 through controlled household testing across common appliances such as televisions, chargers, gaming consoles, and kitchen devices.

\section{Problem Validation and Engineering Design}
\subsection{Customer Needs}
Many households waste electricity through standby loads from televisions, chargers, gaming consoles, printers and kitchen appliances \cite{energygovstandby}. While individual standby loads are small, 
cumulative yearly costs are meaningful. Existing timer switches require manual programming and do not adapt to user behaviour. Most smart plugs offer remote switching but still 
depend on active user involvement. Customers require a solution that automatically reduces unnecessary consumption while maintaining convenience. The system must be simple to use, 
safe, and allow manual override when required.

\subsection{High-Level Objective (HLO-4)}

The subsystem shall automatically recommend or execute outlet shutoff schedules that reduce standby energy consumption by at least 40\% during controlled testing, while allowing manual user override within 2 seconds.

\section{System Requirements}

% \section{System Requirements}

The system requirements for the Automatic Power-off Scheduling subsystem have been developed to ensure alignment with HLO-4 and the broader PowerStake prototype objectives. Requirements are divided into Physical, Functional, and Performance categories. Physical requirements define hardware constraints and design limitations. Functional requirements define what the subsystem must do. Performance requirements define how effectively the subsystem must operate and be validated.

\subsection{Numbering Scheme}

Each requirement is identified using the following format:

\begin{itemize}
    \item REQ-Ph-XX = Physical Requirement
    \item REQ-F-XX = Functional Requirement
    \item REQ-Pe-XX = Performance Requirement
\end{itemize}

\subsection{Physical Requirements}

\begin{table}[H]
\caption{Physical subsystem requirements}
\centering
\renewcommand{\arraystretch}{1.2}
\begin{tabular}{|p{2.5cm}|p{10cm}|p{2cm}|}
\hline
\textbf{ID} & \textbf{Requirement Description} & \textbf{Trace} \\ \hline
REQ-Ph-01 & The subsystem shall control a minimum of four independently switched AC outlets. & HLO-4 \\ \hline
REQ-Ph-02 & The subsystem shall operate from Australian 240V mains supply using isolated switching hardware. & HLO-4 \\ \hline
REQ-Ph-03 & The subsystem shall include wireless communication hardware for dashboard connectivity. & HLO-4 \\ \hline
REQ-Ph-04 & The prototype enclosure footprint shall remain below 350 mm \(\times\) 150 mm. & HLO-4 \\ \hline
REQ-Ph-05 & The subsystem shall utilise low-voltage internal control circuitry separated from switched mains outputs. & HLO-4 \\ \hline
\end{tabular}
\end{table}

\subsection{Functional Requirements}

\begin{table}[H]
\caption{Functional subsystem requirements}
\centering
\renewcommand{\arraystretch}{1.2}
\begin{tabular}{|p{2.5cm}|p{10cm}|p{2cm}|}
\hline
\textbf{ID} & \textbf{Requirement Description} & \textbf{Trace} \\ \hline
REQ-F-01 & The subsystem shall detect outlet power state at intervals no greater than 1 second. & HLO-4 \\ \hline
REQ-F-02 & The subsystem shall identify standby loads below a configurable wattage threshold. & HLO-4 \\ \hline
REQ-F-03 & The subsystem shall generate shutoff recommendations using repeated inactivity patterns. & HLO-4 \\ \hline
REQ-F-04 & The subsystem shall automatically disconnect selected outlets according to approved schedules or inactivity rules. & HLO-4 \\ \hline
REQ-F-05 & The subsystem shall allow manual override through the mobile dashboard. & HLO-4 \\ \hline
REQ-F-06 & The subsystem shall restore outlet power upon authorised user command. & HLO-4 \\ \hline
REQ-F-07 & The subsystem shall log switching actions for later user review. & HLO-4 \\ \hline
\end{tabular}
\end{table}

\subsection{Performance Requirements}

\begin{table}[H]
\caption{Performance subsystem requirements}
\centering
\renewcommand{\arraystretch}{1.2}
\begin{tabular}{|p{2.5cm}|p{10cm}|p{2cm}|}
\hline
\textbf{ID} & \textbf{Requirement Description} & \textbf{Trace} \\ \hline
REQ-Pe-01 & The subsystem shall reduce standby energy consumption by at least 40\% during controlled prototype testing. & HLO-4 \\ \hline
REQ-Pe-02 & Automatic shutoff commands shall execute within 2 seconds of trigger conditions being met. & HLO-4 \\ \hline
REQ-Pe-03 & At least 80\% of schedule recommendations shall be accepted during five representative user test scenarios. & HLO-4 \\ \hline
REQ-Pe-04 & Wireless command success rate shall exceed 95\% during controlled testing. & HLO-4 \\ \hline
REQ-Pe-05 & The subsystem shall operate continuously for 24 hours without requiring reset during bench testing. & HLO-4 \\ \hline
\end{tabular}
\end{table}

\subsection{Verification Matrix}

\begin{table}[H]
\caption{Requirement traceability matrix}
\centering
\renewcommand{\arraystretch}{1.15}
\begin{tabular}{|p{4cm}|c|}
\hline
\textbf{Requirement ID} & \textbf{HLO-4} \\ \hline
REQ-Ph-01 & X \\ \hline
REQ-Ph-02 & X \\ \hline
REQ-Ph-03 & X \\ \hline
REQ-F-01 & X \\ \hline
REQ-F-02 & X \\ \hline
REQ-F-03 & X \\ \hline
REQ-F-04 & X \\ \hline
REQ-F-05 & X \\ \hline
REQ-F-06 & X \\ \hline
REQ-Pe-01 & X \\ \hline
REQ-Pe-02 & X \\ \hline
REQ-Pe-03 & X \\ \hline
REQ-Pe-04 & X \\ \hline
REQ-Pe-05 & X \\ \hline
\end{tabular}
\end{table}

\subsection{Subsystem Architecture}

The AI subsystem forms the analytical layer of the PowerStake platform. It receives historical and near real-time power monitoring data from the onboard circuitry through the Wi-Fi and cloud communication path. Its role is to transform raw electrical measurements into useful outputs such as appliance fault alerts, power-saving opportunities, and recommended automatic shutoff schedules.

The subsystem contains four main stages: data ingestion, preprocessing, inference, and output generation. During data ingestion, measured voltage, current, power, and switching history are received and stored. These data are then filtered, normalised, and converted into useful features such as standby duration, repeated inactivity, and unusual load changes. The inference stage applies analytical logic to classify appliance behaviour as normal, inefficient, or potentially faulty. Finally, the output stage sends recommendations and alerts to the mobile dashboard, where users can review insights and approve or reject suggested actions.

This subsystem interfaces directly with the wider PowerStake monitoring, communication, and dashboard environment, and provides the decision-support capability required for intelligent energy management.

\begin{figure}[H]
    \centering
    \includegraphics[width=\textwidth]{AI_subsystem_architecture.png}
    \caption{Subsystem architecture for the PowerStake AI subsystem \cite{chatgptfigure2026}.}
    \label{fig:ai_subsystem_architecture}
\end{figure}




\subsection{Trade Studies}

A trade study was conducted to determine the most suitable analytical approach for the PowerStake AI subsystem. The subsystem is required to analyse power monitoring data to identify appliance malfunctions, detect power-saving opportunities, and recommend automatic shutoff actions. For the current prototype stage, the preferred solution needed to be accurate enough to demonstrate value, while remaining practical to implement with limited data, modest computational resources, and a short development timeframe.

Three candidate approaches were evaluated: (1) rule-based thresholding, (2) classical machine learning, and (3) lightweight deep learning. These options were compared based on fault detection capability, data requirement, computational demand, interpretability, and prototype suitability.

\begin{table}[H]
\centering
\renewcommand{\arraystretch}{1.15}
\begin{tabular}{|p{2cm}|p{4cm}|p{2cm}|p{3cm}|p{2cm}|}
\hline
\textbf{Option} & \textbf{Description} & \textbf{Data Requirement} & \textbf{Computational Demand} & \textbf{Suitability} \\ \hline
Rule-Based Thresholding & Fixed rules using power, current, and inactivity duration & Low & Low & High \\ \hline
Classical Machine Learning & Feature-based classification using engineered inputs & Medium & Medium & Medium \\ \hline
Lightweight Deep Learning & Learns temporal patterns directly from power traces & High & High & Low \\ \hline
\end{tabular}
\caption{Initial comparison of AI modelling approaches}
\label{tab:ai_options}
\end{table}

\begin{table}[H]
\centering
\renewcommand{\arraystretch}{1.15}
\begin{tabular}{|p{3cm}|c|c|c|c|c|}
\hline
\textbf{Option} & \textbf{Detection} & \textbf{Data Need} & \textbf{Processing} & \textbf{Interpretability} & \textbf{Total} \\ \hline
Rule-Based Thresholding & 3 & 5 & 5 & 5 & 18 \\ \hline
Classical Machine Learning & 4 & 3 & 3 & 3 & 13 \\ \hline
Lightweight Deep Learning & 5 & 2 & 2 & 2 & 11 \\ \hline
\end{tabular}
\caption{Comparison matrix for AI subsystem options}
\label{tab:ai_trade}
\end{table}

Rule-based thresholding was selected as the preferred prototype solution. Although it is less sophisticated than machine learning methods, it is the most 
realistic option for the current project stage because it requires minimal training data, is easy to interpret, and can be implemented rapidly. Classical 
machine learning is a strong future option once more operating data become available. Lightweight deep learning was considered less suitable at this stage 
due to higher computational and data requirements. Overall, rule-based thresholding provides the best balance of practicality, transparency, and prototype 
suitability for demonstrating the AI subsystem value proposition.



\printbibliography


\end{document}

























% ---------------------------------------------------------------------------------














% ---------------------------------------------------------------------
\subsection{old Subsystem Architecture}


Subsystem architecture (Max 200)
Include the block diagram(s) describing the system architecture for the subsystem you oversee.
• Trade studies (300 words max)
Describe the process undertaken for generating concepts and solutions for components of your
subsystem. Include matrices/tables used for decision-making. Select one (1) problems or
components in your system and describe the trade studies undertaken to select a solution.

\textbf{AI SYBSYSTEM }

The Automatic Power-off Scheduling subsystem forms part of the wider PowerStake platform by linking power monitoring, intelligent control, wireless communication, 
and user interaction. Power sensors provide real-time outlet usage data which is processed by the onboard microcontroller. A practical prototype could utilise an 
ESP32 microcontroller due to its integrated Wi-Fi capability, sufficient processing performance, and low-cost implementation \cite{esp32productpage}. As well as ACS712 current sensing modules, 
and isolated relay switching hardware to provide low-cost monitoring and automated outlet control \cite{acs712}.

The controller stores recent usage history and applies scheduling logic to identify repeated inactivity periods or standby behaviour. When suitable conditions are 
met, relay drivers disconnect power to selected outlets in order to reduce unnecessary energy consumption.

The Wi-Fi module enables two-way communication between the device and the mobile dashboard. This allows users to manually override schedules, restore outlet power, 
and review system recommendations. In later development stages, cloud-based analysis could be introduced to improve shutoff recommendations using longer-term usage trends.

\begin{figure}[H]
    \centering
    \includegraphics[width=\textwidth]{sysarch2.png}
    \caption{System Architecture of the Automatic Power-off Scheduling Subsystem \cite{chatgptfigure2026}.}
    \label{fig:sysarch}
\end{figure}
This subsystem interfaces directly with the wider PowerStake monitoring, control, and mobile dashboard environment.


\subsection{Trade Studies}

A trade study was conducted to determine the most suitable switching solution for the Automatic Power-off Scheduling subsystem. This subsystem requires a device capable of disconnecting power to selected outlets once standby conditions or scheduled shutoff periods are detected. The preferred solution needed to provide safe operation, low implementation cost, simple integration with the control hardware, and suitability for prototype development within the available project timeframe.

Three switching approaches were evaluated: (1) electromechanical relay module, (2) solid state relay (SSR), and (3) modified commercial smart plug hardware. Initial comparison of the available options is provided in Table~\ref{tab:optionsummary}.

\begin{table}[H]
    \caption{Initial comparison of switching options}
\centering
\renewcommand{\arraystretch}{1.2}
\begin{tabular}{|p{3.7cm}|p{3.7cm}|p{2cm}|p{3.0cm}|p{2cm}|}
\hline
\textbf{Option} & \textbf{Switching Capability} & \textbf{Cost (AUD)} & \textbf{Integration Complexity} & \textbf{Suitability} \\ \hline
Mechanical Relay Module & 10A, 250VAC switching contacts \cite{songlerelay} & \$8--15 & Low -- common modules available & High \\ \hline
Solid State Relay (SSR) & Silent AC switching, high cycle life \cite{omronssr} & \$15--30 & Medium -- thermal/load matching required & Medium \\ \hline
Modified Smart Plug Hardware & Existing mains switching in consumer enclosure \cite{tplinkplug} & \$20+ & High -- proprietary PCB rework required & Low \\ \hline
\end{tabular}

\label{tab:optionsummary}
\end{table}

The options were then assessed using the criteria of cost, safety, ease of integration, and prototype suitability. Weighting factors were assigned from 1 (lowest importance) to 4 (highest importance), where greater weighting reflects stronger relevance to subsystem requirements.

\begin{table}[H]
    \caption{Trade study weighting factors}
\centering
\renewcommand{\arraystretch}{1.15}
\begin{tabular}{|p{5cm}|c|}
\hline
\textbf{Criterion} & \textbf{Weight} \\ \hline
Cost & 3 \\ \hline
Safety & 4 \\ \hline
Ease of Integration & 4 \\ \hline
Prototype Suitability & 3 \\ \hline
\end{tabular}

\label{tab:weights}
\end{table}

Each option was scored on a scale of 1 to 5, where 5 represents the most favourable outcome. Weighted totals were calculated by multiplying each score by the corresponding criterion weighting.

\begin{table}[H]
    \caption{Weighted comparison matrix for switching solutions}
\centering
\renewcommand{\arraystretch}{1.15}
\begin{tabular}{|p{4.2cm}|c|c|c|c|c|}
\hline
\textbf{Option} & \textbf{Cost} & \textbf{Safety} & \textbf{Integration} & \textbf{Suitability} & \textbf{Total} \\ \hline
Mechanical Relay Module & 5 & 4 & 5 & 5 & 56 \\ \hline
Solid State Relay & 3 & 4 & 3 & 3 & 41 \\ \hline
Modified Smart Plug Hardware & 2 & 2 & 1 & 2 & 21 \\ \hline
\end{tabular}

\label{tab:trade}
\end{table}

The mechanical relay module achieved the highest total score and was selected as the preferred solution. Relay boards are inexpensive, widely available, and provide galvanic isolation between the low-voltage control circuit and the switched AC load, making them appropriate for prototype testing and early subsystem validation \cite{songlerelay}.

Solid state relays offer silent operation and long switching life; however, they are more expensive and may require additional thermal management depending on the connected load \cite{omronssr}. While technically effective, they were considered less practical for the current prototype stage.

Modified smart plug hardware ranked lowest due to proprietary circuitry, difficult modification, and reduced transparency during testing. Although these devices contain integrated switching hardware, they are generally not designed for redevelopment or subsystem prototyping \cite{tplinkplug}.

Overall, the mechanical relay module provided the best balance of cost, safety, ease of implementation, and suitability for demonstrating the subsystem High-Level Objective within the current prototype scope.


















\subsection{Trade Studies}

A trade study was completed to determine the most suitable outlet switching technology for the Automatic Power-off Scheduling subsystem. This component is required 
to disconnect power to selected outlets when standby conditions or scheduled shutoff periods are detected. As the project is currently at concept stage and aims to
 reach TRL 3 progressing toward early TRL 4, practical factors such as low cost, simple integration, safe operation, and rapid prototyping were prioritised over 
 premium commercial performance.

Three switching options were considered: (1) electromechanical relay module, (2) solid state relay (SSR), and (3) modified commercial smart plug hardware. These 
options were compared using the criteria shown in Table~\ref{tab:weights}.

\begin{table}[H]
\centering
\renewcommand{\arraystretch}{1.15}
\begin{tabular}{|p{5cm}|c|}
\hline
\textbf{Criterion} & \textbf{Weight (\%)} \\ \hline
Electrical Safety / Isolation & 30 \\ \hline
Cost & 25 \\ \hline
Ease of Integration & 20 \\ \hline
Reliability & 15 \\ \hline
Prototype Delivery Speed & 10 \\ \hline
\end{tabular}
\caption{Trade study weighting criteria}
\label{tab:weights}
\end{table}

Each option was scored from 1 (poor) to 5 (excellent). Final weighted results are presented in Table~\ref{tab:trade}.

\begin{table}[H]
\centering
\renewcommand{\arraystretch}{1.15}
\begin{tabular}{|p{6cm}|c|c|}
\hline
\textbf{Option} & \textbf{Weighted Score} & \textbf{Rank} \\ \hline
Mechanical Relay Module & 4.45 & 1 \\ \hline
Solid State Relay & 3.85 & 2 \\ \hline
Modified Smart Plug Hardware & 2.95 & 3 \\ \hline
\end{tabular}
\caption{Weighted comparison of switching technologies}
\label{tab:trade}
\end{table}

The mechanical relay module achieved the highest score and will be selected for the prototype build. Relay modules are inexpensive, commonly available, and can be 
directly controlled using ESP32 digital outputs \cite{esp32datasheet}. They also provide electrical isolation between the low-voltage control circuit and higher-voltage switching side, 
which is beneficial during early prototype development.

Solid state relays offer quieter operation and longer switching life, however they are more expensive and may require additional thermal considerations. Modified 
smart plug hardware was ranked lowest due to limited accessibility, proprietary circuitry, and added integration difficulty.

Overall, the relay module was considered the most practical and realistic option for demonstrating subsystem functionality within the available unit timeframe.


% ---------------------------- v1 

\subsection{Trade Studies}

A trade study was conducted to determine the most suitable switching method for the Automatic Power-off Scheduling subsystem. Reliable switching is essential, as the selected component must 
disconnect outlet power when standby conditions or scheduled shutoff periods are detected. For the current prototype stage, the preferred solution needed to be low-cost, safe, easy to integrate 
with an ESP32 microcontroller, and practical to implement within the available project timeframe.

Three switching approaches were evaluated: (1) electromechanical relay module, (2) solid state relay (SSR), and (3) modified commercial smart plug hardware. These options were compared based 
on switching capability, cost, integration complexity, prototype suitability, and overall practicality.

\begin{table}[H]
\centering
\renewcommand{\arraystretch}{1.2}
\begin{tabular}{|p{3.5cm}|p{3.8cm}|p{2cm}|p{3.2cm}|p{2cm}|}
\hline
\textbf{Option} & \textbf{Switching Capability} & \textbf{Cost (AUD)} & \textbf{Integration Complexity} & \textbf{Suitability} \\ \hline
Mechanical Relay Module & 10A, 250VAC switching contacts \cite{songlerelay} & \$8--15 & Low -- direct GPIO control, common modules available & High \\ \hline
Solid State Relay (SSR) & Silent AC switching, high cycle life \cite{omronssr} & \$15--30 & Medium -- load matching and thermal management required & Medium \\ \hline
Modified Smart Plug Hardware & Existing mains switching in consumer enclosure \cite{tplinkplug} & \$20+ & High -- proprietary PCB, difficult modification & Low \\ \hline
\end{tabular}
\caption{Automatic Power-off Scheduling trade study}
\label{tab:trade}
\end{table}

The mechanical relay module was selected as the preferred solution due to its low cost, simple control interface, and suitability for rapid prototype development. Standard relay modules can 
be controlled directly using ESP32 digital outputs and provide galvanic isolation between the low-voltage control circuit and the switched load \cite{songlerelay}. This makes them practical for early-stage testing and subsystem validation.

Solid state relays provide silent operation and long switching life; however, they are more expensive and may require additional thermal management depending on the connected 
load \cite{omronssr}. While technically strong, they were considered less practical for an initial university prototype.

Modified smart plug hardware was the least suitable option. Although it contains integrated switching hardware, internal circuitry is often proprietary and not intended for modification, 
increasing development time and reducing transparency during testing \cite{tplinkplug}.

Overall, the mechanical relay module offered the best balance of cost, safety, ease of integration, and functional performance for achieving the subsystem High-Level Objective within the 
current project scope.
\newpage
\printbibliography



% ---------------------------------------------------------------------------------



\subsection{Physical Requirements}

\begin{table}[H]
\caption{Physical system requirements}
\centering
\renewcommand{\arraystretch}{1.2}
\begin{tabular}{|p{2cm}|p{10.5cm}|p{2cm}|}
\hline
\textbf{ID} & \textbf{Requirement} & \textbf{HLO} \\ \hline
PH-01 & Subsystem shall control a minimum of four independent outlets. & HLO-4 \\ \hline
PH-02 & Subsystem shall operate from 240V Australian mains supply via an isolated relay module. & HLO-4 \\ \hline
PH-03 & Subsystem shall include a Wi-Fi capable microcontroller for communication and control. & HLO-4 \\ \hline
PH-04 & Prototype enclosure footprint shall be less than 350 mm \(\times\) 150 mm. & HLO-4 \\ \hline
\end{tabular}

\end{table}

\subsection{Functional Requirements}

\begin{table}[H]
    \caption{Functional system requirements}
\centering
\renewcommand{\arraystretch}{1.2}
\begin{tabular}{|p{2cm}|p{10.5cm}|p{2cm}|}
\hline
\textbf{ID} & \textbf{Requirement} & \textbf{HLO} \\ \hline
FN-01 & Subsystem shall measure outlet power state once per second. & HLO-4 \\ \hline
FN-02 & Subsystem shall detect standby condition below a configurable wattage threshold. & HLO-4 \\ \hline
FN-03 & Subsystem shall generate shutoff schedules using repeated inactivity patterns. & HLO-4 \\ \hline
FN-04 & Subsystem shall allow manual override through the mobile dashboard. & HLO-4 \\ \hline
FN-05 & Subsystem shall restore power immediately upon authorised user command. & HLO-4 \\ \hline
\end{tabular}

\end{table}

\subsection{Performance Requirements}

\begin{table}[H]
    \caption{Performance system requirements}
\centering
\renewcommand{\arraystretch}{1.2}
\begin{tabular}{|p{2cm}|p{10.5cm}|p{2cm}|}
\hline
\textbf{ID} & \textbf{Requirement} & \textbf{HLO} \\ \hline
PE-01 & Subsystem shall reduce standby energy consumption by at least 40\% during controlled testing. & HLO-4 \\ \hline
PE-02 & Switching response time shall be less than 2 seconds after command input. & HLO-4 \\ \hline
PE-03 & Schedule recommendations shall achieve at least 80\% user acceptance accuracy. & HLO-4 \\ \hline
PE-04 & Subsystem shall operate continuously for 30 days without requiring reset. & HLO-4 \\ \hline
\end{tabular}

\end{table}